\documentclass[12pt,a4paper]{article}

\usepackage[a4paper,margin=2.5cm]{geometry}
\usepackage{setspace}
\usepackage{titlesec}
\usepackage{enumitem}
\usepackage{tabularx}
\usepackage{longtable}
\usepackage{booktabs}
\usepackage{xcolor}
\usepackage{hyperref}
\usepackage{listings}
\usepackage{fancyhdr}
\usepackage[utf8]{inputenc}
\usepackage[T1]{fontenc}
\usepackage{mathptmx}
\usepackage[vietnamese]{babel}
\usepackage{amsmath}

\onehalfspacing

\hypersetup{
  colorlinks=true,
  linkcolor=blue,
  urlcolor=blue
}

% \pagestyle{fancy}
% \fancyhf{}
% \rhead{WebRTC TURN + Room + Group Call}
% \lhead{Report}
% \cfoot{\thepage}

\titleformat{\section}{\large\bfseries}{\thesection.}{0.5em}{}
\titleformat{\subsection}{\normalsize\bfseries}{\thesubsection}{0.5em}{}

\lstdefinelanguage{json}{
  basicstyle=\ttfamily\small,
  morestring=[b]",
  stringstyle=\color{blue},
  showstringspaces=false,
  breaklines=true
}

\lstset{
  frame=single,
  framerule=0.3pt,
  rulecolor=\color{gray},
  backgroundcolor=\color{gray!5},
  breaklines=true,
  columns=fullflexible
}

\begin{document}

\begin{center}
    {\LARGE \textbf{Report WebRTC}}\\[0.5em]
    {\large WebRTC Group Call System with Metered TURN, Room, and Mesh Topology}\\[2em]

    % Phần thông tin Giảng viên và Sinh viên chia cột
    \begin{tabular}{p{0.15\textwidth} l p{0.1\textwidth} p{0.15\textwidth} l}
        \textbf{Lecturer:}   & Lê Viết Long         & & \textbf{Students:} & Trần Bình Duy - 23127xxx \\
        \textbf{} & Nguyễn Thanh Quân     & &                    & Nguyễn Văn Khánh - 23127388 \\

    \end{tabular}
\end{center}

% \begin{center}
%   {\LARGE \textbf{Report\_WebRTC\_TURN\_Room\_GroupCall}}\\[0.5em]
%   {\large WebRTC Group Call System with Metered TURN, Room, and Mesh Topology}\\[1em]
%   Sinh viên: \underline{\hspace{7cm}}\\[0.5em]
%   Môn học: \underline{\hspace{7cm}}\\[0.5em]
%   Ngày nộp: \underline{\hspace{6cm}}
% \end{center}

\vspace{1em}
\hrule
\vspace{1em}

\section{Mô tả kiến trúc hệ thống}

\subsection{Tổng quan client/server}
Hệ thống được xây dựng theo mô hình signaling server + WebRTC peer-to-peer:
\begin{itemize}[leftmargin=1.5em]
  \item \textbf{Server (Node.js + HTTPS + WebSocket):} quản lý đăng ký nickname, quản lý phòng, phát tán danh sách thành viên, chuyển tiếp signaling messages (offer/answer/candidate).
  \item \textbf{Client (HTML/CSS/JS):} lấy media bằng \texttt{getUserMedia()}, thiết lập \texttt{RTCPeerConnection}, hiển thị grid video, xử lý trạng thái cuộc gọi.
  \item \textbf{TURN service (Metered.ca):} cung cấp relay khi kết nối trực tiếp thất bại do NAT/Firewall.
\end{itemize}

\subsection{Luồng signaling và media}
\begin{enumerate}[leftmargin=1.8em]
  \item Client đăng ký vào signaling server qua WebSocket.
  \item Người dùng tạo phòng hoặc tham gia phòng.
  \item Khi bắt đầu gọi nhóm, các peer trao đổi SDP offer/answer và ICE candidates qua server.
  \item Sau khi thương lượng thành công, media truyền trực tiếp P2P nếu khả dụng; nếu không, fallback qua TURN relay.
\end{enumerate}

\subsection{WebRTC trong hệ thống này}
WebRTC là lớp truyền thông thời gian thực giữa các trình duyệt, chịu trách nhiệm vận chuyển audio/video sau khi signaling hoàn tất.

\textbf{Các thành phần WebRTC được sử dụng:}
\begin{itemize}[leftmargin=1.5em]
  \item \texttt{getUserMedia()}: lấy luồng camera/microphone từ thiết bị người dùng.
  \item \texttt{RTCPeerConnection}: tạo và duy trì kết nối media giữa 2 peers.
  \item \texttt{MediaStreamTrack}: quản lý track audio/video để bật/tắt mic/camera.
  \item \texttt{getStats()}: thu thập thống kê chất lượng kết nối (RTT, packet loss, FPS, candidate type).
\end{itemize}

\textbf{Vòng đời kết nối WebRTC trong app:}
\begin{enumerate}[leftmargin=1.8em]
  \item Tạo \texttt{RTCPeerConnection} với danh sách \texttt{iceServers} do server cung cấp.
  \item Thêm local tracks (audio/video) vào peer connection.
  \item Trao đổi SDP offer/answer qua signaling server.
  \item Trao đổi ICE candidates để tìm đường truyền tối ưu.
  \item Vào trạng thái \texttt{connected}: nhận và hiển thị remote stream trên video grid.
  \item Khi mạng gián đoạn: theo dõi \texttt{iceConnectionState}, thực hiện ICE restart nếu cần.
\end{enumerate}

\textbf{Vai trò của WebRTC trong kiến trúc:}
\begin{itemize}[leftmargin=1.5em]
  \item Signaling server chỉ điều phối tín hiệu, không truyền media.
  \item Media đi trực tiếp giữa peers (P2P) hoặc đi qua TURN relay khi P2P thất bại.
  \item Với group call mesh, mỗi user duy trì nhiều \texttt{RTCPeerConnection} song song (mỗi peer một kết nối).
\end{itemize}

\subsection{Mô hình ICE/STUN/TURN}
\textbf{ICE (Interactive Connectivity Establishment)} là cơ chế để hai peer tìm đường truyền media tốt nhất giữa nhiều khả năng kết nối.

\textbf{Các loại ICE candidate:}
\begin{itemize}[leftmargin=1.5em]
  \item \texttt{host}: candidate nội mạng (LAN/direct), độ trễ thấp nhất khi hai thiết bị có thể kết nối trực tiếp.
  \item \texttt{srflx}: candidate phản xạ từ STUN (public IP/port mapping), dùng để vượt NAT thông thường.
  \item \texttt{relay}: candidate đi qua TURN relay (Metered), dùng khi \texttt{host}/\texttt{srflx} không tạo được đường truyền end-to-end.
\end{itemize}

\textbf{Quy trình ICE trong hệ thống:}
\begin{enumerate}[leftmargin=1.8em]
  \item Mỗi peer thu thập candidate từ danh sách \texttt{iceServers} (STUN + TURN).
  \item Hai bên trao đổi candidate qua signaling server bằng message \texttt{candidate}.
  \item WebRTC thực hiện connectivity checks trên các candidate pairs khả dụng.
  \item Chọn candidate pair tốt nhất để truyền media; trạng thái chuyển từ \texttt{checking} sang \texttt{connected/completed}.
\end{enumerate}

\textbf{Thứ tự ưu tiên và fallback:}
\begin{itemize}[leftmargin=1.5em]
  \item Thông thường ưu tiên đường direct (\texttt{host} hoặc \texttt{srflx}) để giảm độ trễ và tiết kiệm relay bandwidth.
  \item Nếu direct path thất bại (NAT đối xứng, firewall chặn UDP/TCP, mạng di động hạn chế), hệ thống fallback sang \texttt{relay (TURN)}.
  \item Trong ứng dụng này, khi \texttt{iceConnectionState} chuyển sang \texttt{disconnected/failed}, client kích hoạt ICE restart để thử tái thiết lập đường truyền.
\end{itemize}

\textbf{Ý nghĩa trong kiểm thử:}
\begin{itemize}[leftmargin=1.5em]
  \item Test LAN thường thấy \texttt{host} $\Rightarrow$ P2P trực tiếp.
  \item Test khác mạng (WiFi vs 4G) thường xuất hiện \texttt{relay} $\Rightarrow$ TURN hoạt động đúng.
  \item Thống kê từ \texttt{getStats()} được dùng để xác nhận candidate type cuối cùng và chứng minh relay trong báo cáo.
\end{itemize}

\section{Định nghĩa protocol signaling}

\subsection{Danh sách message JSON chính}
\begin{longtable}{@{}p{3cm}p{2cm}p{8.5cm}@{}}
\toprule
\textbf{Message} & \textbf{Hướng} & \textbf{Payload} \\
\midrule
\endhead
\texttt{register} & C$\rightarrow$S & \texttt{\{type, name\}} \\
\texttt{registered} & S$\rightarrow$C & \texttt{\{type, name, iceServers[]\}} \\
\texttt{createRoom} & C$\rightarrow$S & \texttt{\{type, roomId, name\}} \\
\texttt{joinRoom} & C$\rightarrow$S & \texttt{\{type, roomId, name\}} \\
\texttt{roomJoined} & S$\rightarrow$C & \texttt{\{type, roomId, callActive\}} \\
\texttt{roomMembers} & S$\rightarrow$C & \texttt{\{type, roomId, members[], callActive\}} \\
\texttt{offer} & C$\rightarrow$S$\rightarrow$C & \texttt{\{type, roomId, sender, target, offer\}} \\
\texttt{answer} & C$\rightarrow$S$\rightarrow$C & \texttt{\{type, roomId, sender, target, answer\}} \\
\texttt{candidate} & C$\rightarrow$S$\rightarrow$C & \texttt{\{type, roomId, sender, target, candidate\}} \\
\texttt{leaveRoom} & C$\rightarrow$S & \texttt{\{type, roomId, sender\}} \\
\texttt{memberLeft} & S$\rightarrow$C & \texttt{\{type, roomId, name\}} \\
\texttt{endCall} & C$\rightarrow$S & \texttt{\{type, roomId, sender\}} \\
\texttt{peerEndedCall} & S$\rightarrow$C & \texttt{\{type, roomId, name\}} \\
\texttt{error} & S$\rightarrow$C & \texttt{\{type, message\}} \\
\bottomrule
\end{longtable}

\subsection{Ví dụ gói tin}
\begin{lstlisting}[language=json,caption={Offer message}]
{
  "type": "offer",
  "roomId": "abc123",
  "sender": "userA",
  "target": "userB",
  "offer": {
    "type": "offer",
    "sdp": "v=0..."
  }
}
\end{lstlisting}

\begin{lstlisting}[language=json,caption={ICE candidate message}]
{
  "type": "candidate",
  "roomId": "abc123",
  "sender": "userA",
  "target": "userB",
  "candidate": {
    "candidate": "candidate:...",
    "sdpMid": "0",
    "sdpMLineIndex": 0
  }
}
\end{lstlisting}

\section{Thiết kế room và group call}

\subsection{Quản lý room}
Server quản lý cấu trúc:
\begin{itemize}[leftmargin=1.5em]
  \item \texttt{clients}: ánh xạ \texttt{name $\rightarrow$ \{ws, roomId\}}
  \item \texttt{rooms}: ánh xạ \texttt{roomId $\rightarrow$ \{members: Set, callActive\}}
\end{itemize}

Khi có user join/leave, server phát \texttt{roomMembers} để cập nhật realtime cho toàn bộ thành viên trong phòng.

\subsection{Group call mesh topology}
Với phòng có $n$ người, mỗi client tạo $n-1$ peer connections:
\[
\text{Tổng số kết nối} = \frac{n(n-1)}{2}
\]

\textbf{Ví dụ:} 4 người $\Rightarrow$ 6 kết nối song song.

\subsection{Hiển thị grid video}
\begin{itemize}[leftmargin=1.5em]
  \item 1 local video.
  \item Nhiều remote video (mỗi peer một khung).
  \item Grid tự điều chỉnh theo số participants (1, 2, 3--4, 5--6).
  \item Khi peer rời cuộc gọi/phòng: đóng PC tương ứng và remove video element để tránh khung bị đứng hình.
\end{itemize}

\section{Triển khai TURN (dịch vụ Metered.ca)}

\subsection{Cấu hình sử dụng}
Hệ thống sử dụng dịch vụ TURN của Metered.ca. Cấu hình trong file \texttt{.env}:
\begin{lstlisting}[language=json,caption={Environment configuration}]
PORT=3000
METERED_APP_NAME=your_app_name
METERED_API_KEY=your_api_key
\end{lstlisting}

\subsection{Ports và transport}
Khi truy vấn endpoint credentials của Metered:
\begin{lstlisting}
https://webrtclab.metered.live/api/v1/turn/credentials?apiKey=YOUR_API_KEY
\end{lstlisting}

API trả về danh sách ICE endpoints thực tế như sau:
\begin{lstlisting}[language=json,caption={Metered ICE endpoints (actual response format)}]
[
  {"urls":"stun:stun.relay.metered.ca:80"},
  {"urls":"turn:standard.relay.metered.ca:80","username":"<USERNAME>","credential":"<CREDENTIAL>"},
  {"urls":"turn:standard.relay.metered.ca:80?transport=tcp","username":"<USERNAME>","credential":"<CREDENTIAL>"},
  {"urls":"turn:standard.relay.metered.ca:443","username":"<USERNAME>","credential":"<CREDENTIAL>"},
  {"urls":"turns:standard.relay.metered.ca:443?transport=tcp","username":"<USERNAME>","credential":"<CREDENTIAL>"}
]
\end{lstlisting}

Từ response trên, các port/transport sử dụng gồm:
\begin{itemize}[leftmargin=1.5em]
  \item \texttt{stun:stun.relay.metered.ca:80} (STUN).
  \item \texttt{turn:standard.relay.metered.ca:80} (TURN mặc định, thường ưu tiên UDP).
  \item \texttt{turn:standard.relay.metered.ca:80?transport=tcp} (TURN over TCP).
  \item \texttt{turn:standard.relay.metered.ca:443} (TURN qua port 443).
  \item \texttt{turns:standard.relay.metered.ca:443?transport=tcp} (TURN over TLS).
\end{itemize}


\subsection{Credential}
Credential được cấp động từ API của Metered dựa trên:
\begin{itemize}[leftmargin=1.5em]
  \item App Name
  \item API Key
\end{itemize}
Server cache danh sách credentials để giảm số lần gọi API.

\section{Kết quả kiểm thử}

\subsection{Test LAN (P2P)}
\textbf{Môi trường:} nhiều client cùng mạng LAN/WiFi.

\textbf{Kết quả:}
\begin{itemize}[leftmargin=1.5em]
  \item Gọi nhóm hoạt động ổn định.
  \item Candidate chủ yếu là \texttt{host}.
  \item \texttt{connectionState}: \texttt{connected}
  \item \texttt{iceConnectionState}: \texttt{connected}
\end{itemize}

\textbf{Log mẫu:}
\begin{lstlisting}
[10:57:47] PC [test2] iceConnectionState: connected
[10:57:47] PC [test2] connectionState: connected
[10:57:47] [test2] Candidate: local=host, remote=host -> host (LAN)
\end{lstlisting}

\subsection{Test khác mạng / 4G (TURN relay)}
\textbf{Môi trường:} 1 máy WiFi + 1 máy 4G (khác NAT).

\textbf{Kết quả:}
\begin{itemize}[leftmargin=1.5em]
  \item ICE gathering thu được \texttt{host}, \texttt{srflx}, và \texttt{relay}.
  \item Kết nối cuối cùng chọn \texttt{relay (TURN)} khi direct path không khả dụng.
  \item \texttt{connectionState}: \texttt{connected}
  \item \texttt{iceConnectionState}: \texttt{connected}
\end{itemize}

\textbf{Log mẫu (thực tế):}
\begin{lstlisting}
[13:17:36] ICE candidate [Dt]: host
[13:17:36] ICE candidate [Dt]: srflx (STUN)
[13:17:37] ICE candidate [Dt]: relay (TURN)
[13:17:37] PC [Dt] iceConnectionState: connected
[13:17:38] PC [Dt] connectionState: connected
[13:17:38] [Dt] Candidate: local=relay, remote=relay -> relay (TURN)
\end{lstlisting}

\subsection{Thống kê từ getStats()}
Ứng dụng ghi nhận:
\begin{itemize}[leftmargin=1.5em]
  \item loại candidate pair đang dùng (host/srflx/relay),
  \item RTT (\texttt{currentRoundTripTime}),
  \item FPS video inbound,
  \item packet loss.
\end{itemize}

\section{Hạn chế và hướng phát triển}

\subsection{Hạn chế hiện tại}
\begin{itemize}[leftmargin=1.5em]
  \item Mesh topology tiêu tốn băng thông theo $O(n^2)$, khó scale cho phòng lớn.
  \item Chưa có SFU/MCU để tối ưu media routing.
  \item Dùng self-signed cert trong môi trường local gây thêm bước trust certificate.
\end{itemize}

\subsection{Hướng phát triển}
\begin{itemize}[leftmargin=1.5em]
  \item Chuyển sang SFU (mediasoup/Janus) để hỗ trợ nhiều người hơn.
  \item Adaptive bitrate và simulcast để tối ưu chất lượng theo mạng.
  \item Cải thiện UI/UX: layout linh hoạt hơn, trạng thái mạng trực quan hơn.
  \item Bổ sung screen sharing, chat, recording.
\end{itemize}


\end{document}

